%% othersock.tex
%%
%% Copyright (C) 2004-2018 Simone Piccardi.  Permission is granted to
%% copy, distribute and/or modify this document under the terms of the GNU Free
%% Documentation License, Version 1.1 or any later version published by the
%% Free Software Foundation; with the Invariant Sections being "Un preambolo",
%% with no Front-Cover Texts, and with no Back-Cover Texts.  A copy of the
%% license is included in the section entitled "GNU Free Documentation
%% License".
%%

\chapter{Gli altri tipi di socket}
\label{cha:other_socket}

Dopo aver trattato in cap.~\ref{cha:TCP_socket} i socket TCP, che
costituiscono l'esempio più comune dell'interfaccia dei socket, esamineremo in
questo capitolo gli altri tipi di socket, a partire dai socket UDP, e i socket
\textit{Unix domain} già incontrati in sez.~\ref{sec:ipc_socketpair}.


\section{I socket UDP}
\label{sec:UDP_socket}

Dopo i socket TCP i socket più utilizzati nella programmazione di rete sono i
socket UDP: protocolli diffusi come NFS o il DNS usano principalmente questo
tipo di socket. Tratteremo in questa sezione le loro caratteristiche
principali e le modalità per il loro utilizzo.


\subsection{Le caratteristiche di un socket UDP}
\label{sec:UDP_characteristics}

Come illustrato in sez.\ref{sec:net_udp} UDP è un protocollo molto semplice che
non supporta le connessioni e non è affidabile: esso si appoggia direttamente
sopra IP (per i dettagli sul protocollo si veda sez.~\ref{sec:udp_protocol}).
I dati vengono inviati in forma di pacchetti, e non ne è assicurata né la
effettiva ricezione né l'arrivo nell'ordine in cui vengono inviati. Il
vantaggio del protocollo è la velocità, non è necessario trasmettere le
informazioni di controllo ed il risultato è una trasmissione di dati più
veloce ed immediata.

Questo significa che a differenza dei socket TCP i socket UDP non supportano
una comunicazione di tipo \textit{stream} in cui si ha a disposizione un
flusso continuo di dati che può essere letto un po' alla volta, ma piuttosto
una comunicazione di tipo \textit{datagram}, in cui i dati arrivano in singoli
blocchi che devono essere letti integralmente.

Questo diverso comportamento significa anche che i socket UDP, pur
appartenendo alla famiglia \const{PF\_INET}\footnote{o \const{PF\_INET6}
  qualora si usasse invece il protocollo IPv6, che pure supporta UDP.} devono
essere aperti quando si usa la funzione \func{socket} (si riveda quanto
illustrato a suo tempo in tab.~\ref{tab:sock_sock_valid_combinations})
utilizzando per il tipo di socket il valore \const{SOCK\_DGRAM}.

Questa differenza comporta ovviamente che anche le modalità con cui si usano i
socket UDP sono completamente diverse rispetto ai socket TCP, ed in
particolare non esistendo il concetto di connessione non esiste il meccanismo
del \textit{three way handshake} né quello degli stati del protocollo. In
realtà tutto quello che avviene nella comunicazione attraverso dei socket UDP
è la trasmissione di un pacchetto da un client ad un server o viceversa,
secondo lo schema illustrato in fig.~\ref{fig:UDP_packet-exchange}.

\begin{figure}[!htb]
  \centering \includegraphics[width=10cm]{img/udp_connection}  
  \caption{Lo schema di interscambio dei pacchetti per una comunicazione via
     UDP.}
  \label{fig:UDP_packet-exchange}
\end{figure}

Come illustrato in fig.~\ref{fig:UDP_packet-exchange} la struttura generica di
un server UDP prevede, una volta creato il socket, la chiamata a \func{bind}
per mettersi in ascolto dei dati, questa è l'unica parte comune con un server
TCP. Non essendovi il concetto di connessione le funzioni \func{listen} ed
\func{accept} non sono mai utilizzate nel caso di server UDP. La ricezione dei
dati dal client avviene attraverso la funzione \func{recvfrom}, mentre una
eventuale risposta sarà inviata con la funzione \func{sendto}.

Da parte del client invece, una volta creato il socket non sarà necessario
connettersi con \func{connect} (anche se, come vedremo in
sez.~\ref{sec:UDP_connect}, è possibile usare questa funzione, con un
significato comunque diverso) ma si potrà effettuare direttamente una
richiesta inviando un pacchetto con la funzione \func{sendto} e si potrà
leggere una eventuale risposta con la funzione \func{recvfrom}.

Anche se UDP è completamente diverso rispetto a TCP resta identica la
possibilità di gestire più canali di comunicazione fra due macchine
utilizzando le porte. In questo caso il server dovrà usare comunque la
funzione \func{bind} per scegliere la porta su cui ricevere i dati, e come nel
caso dei socket TCP si potrà usare il comando \cmd{netstat} per
verificare quali socket sono in ascolto:
\begin{verbatim}
[piccardi@gont gapil]# netstat -anu
Active Internet connections (servers and established)
Proto Recv-Q Send-Q Local Address           Foreign Address         State
udp        0      0 0.0.0.0:32768           0.0.0.0:*
udp        0      0 192.168.1.2:53          0.0.0.0:*
udp        0      0 127.0.0.1:53            0.0.0.0:*
udp        0      0 0.0.0.0:67              0.0.0.0:*
\end{verbatim}
in questo caso abbiamo attivi il DNS (sulla porta 53, e sulla 32768 per la
connessione di controllo del server \cmd{named}) ed un server DHCP (sulla
porta 67).

Si noti però come in questo caso la colonna che indica lo stato sia vuota. I
socket UDP infatti non hanno uno stato. Inoltre anche in presenza di traffico
non si avranno indicazioni delle connessioni attive, proprio perché questo
concetto non esiste per i socket UDP, il kernel si limita infatti a ricevere i
pacchetti ed inviarli al processo in ascolto sulla porta cui essi sono
destinati, oppure a scartarli inviando un messaggio \textit{ICMP port
  unreachable} qualora non vi sia nessun processo in ascolto.


\subsection{Le funzioni \func{sendto} e \func{recvfrom}}
\label{sec:UDP_sendto_recvfrom}

Come accennato in sez.~\ref{sec:UDP_characteristics} le due funzioni
principali usate per la trasmissione di dati attraverso i socket UDP (ed in
generale per i socket di tipo \textit{datagram}) sono \func{sendto} e
\func{recvfrom}. La necessità di usare queste funzioni è dovuta al fatto che
non esistendo con UDP il concetto di connessione, non si può stabilire (come
avviene con i socket TCP grazie alla chiamata ad \func{accept} che li associa
ad una connessione) quali sono la sorgente e la destinazione dei dati che
passano sul socket.\footnote{anche se in alcuni casi, come quello di un client
  che contatta un server, è possibile connettere il socket (vedi
  sez.~\ref{sec:UDP_connect}), la necessità di \func{sendto} e \func{recvfrom}
  resta, dato che questo è possibile solo sul lato del client.}

Per questo anche se in generale si possono comunque leggere i dati con
\func{read}, usando questa funzione non si sarà in grado di determinare da
quale fra i possibili corrispondenti (se ve ne sono più di uno, come avviene
sul lato del server) questi arrivino. E non sarà comunque possibile usare
\func{write} (che fallisce un errore di \errval{EDESTADDRREQ}) in quanto non
è determinato la destinazione che i dati avrebbero.

Per questo motivo nel caso di UDP diventa essenziale utilizzare queste due
funzioni, che sono comunque utilizzabili in generale per la trasmissione di
dati attraverso qualunque tipo di socket. Esse hanno la caratteristica di
prevedere tre argomenti aggiuntivi attraverso i quali è possibile specificare
la destinazione dei dati trasmessi o ottenere l'origine dei dati ricevuti. La
prima di queste funzioni è \funcd{sendto} ed il suo prototipo\footnote{il
  prototipo illustrato è quello utilizzato dalla \acr{glibc}, che segue le
  \textit{Single Unix Specification}, l'argomento \param{flags} era di tipo
  \ctyp{int} nei vari BSD4.*, mentre nelle \acr{libc4} e \acr{libc5} veniva
  usato un \texttt{unsigned int}; l'argomento \param{len} era \ctyp{int} nei
  vari BSD4.* e nella \acr{libc4}, ma \type{size\_t} nella \acr{libc5}; infine
  l'argomento \param{tolen} era \ctyp{int} nei vari BSD4.*, nella \acr{libc4} e
  nella \acr{libc5}.} è:
\begin{functions}
  \headdecl{sys/types.h}
  \headdecl{sys/socket.h}
  
  \funcdecl{ssize\_t sendto(int sockfd, const void *buf, size\_t len, int
    flags, const struct sockaddr *to, socklen\_t tolen)}
  
  Trasmette un messaggio ad un altro socket.
  
  \bodydesc{La funzione restituisce il numero di caratteri inviati in caso di
    successo e -1 per un errore; nel qual caso \var{errno} viene impostata al
    rispettivo codice di errore:
  \begin{errlist}
  \item[\errcode{EAGAIN}] il socket è in modalità non bloccante, ma
    l'operazione richiede che la funzione si blocchi.
  \item[\errcode{ECONNRESET}] l'altro capo della comunicazione ha resettato la
    connessione.
  \item[\errcode{EDESTADDRREQ}] il socket non è di tipo connesso, e non si è
    specificato un indirizzo di destinazione.
  \item[\errcode{EISCONN}] il socket è già connesso, ma si è specificato un
    destinatario.
  \item[\errcode{EMSGSIZE}] il tipo di socket richiede l'invio dei dati in un
    blocco unico, ma la dimensione del messaggio lo rende impossibile.
  \item[\errcode{ENOBUFS}] la coda di uscita dell'interfaccia è già piena (di
    norma Linux non usa questo messaggio ma scarta silenziosamente i
    pacchetti).
  \item[\errcode{ENOTCONN}] il socket non è connesso e non si è specificata
    una destinazione.
  \item[\errcode{EOPNOTSUPP}] il valore di \param{flag} non è appropriato per
    il tipo di socket usato.
  \item[\errcode{EPIPE}] il capo locale della connessione è stato chiuso, si
    riceverà anche un segnale di \signal{SIGPIPE}, a meno di non aver impostato
    \const{MSG\_NOSIGNAL} in \param{flags}.
  \end{errlist}
  ed anche \errval{EFAULT}, \errval{EBADF}, \errval{EINVAL}, \errval{EINTR},
  \errval{ENOMEM}, \errval{ENOTSOCK} più gli eventuali altri errori relativi
  ai protocolli utilizzati.}
\end{functions}

I primi tre argomenti sono identici a quelli della funzione \func{write} e
specificano il socket \param{sockfd} a cui si fa riferimento, il buffer
\param{buf} che contiene i dati da inviare e la relativa lunghezza
\param{len}.  Come per \func{write} la funzione ritorna il numero di byte
inviati; nel caso di UDP però questo deve sempre corrispondere alla dimensione
totale specificata da \param{len} in quanto i dati vengono sempre inviati in
forma di pacchetto e non possono essere spezzati in invii successivi.  Qualora
non ci sia spazio nel buffer di uscita la funzione si blocca (a meno di non
avere aperto il socket in modalità non bloccante), se invece non è possibile
inviare il messaggio all'interno di un unico pacchetto (ad esempio perché
eccede le dimensioni massime del protocollo sottostante utilizzato) essa
fallisce con l'errore di \errcode{EMSGSIZE}. 

I due argomenti \param{to} e \param{tolen} servono a specificare la
destinazione del messaggio da inviare, e indicano rispettivamente la struttura
contenente l'indirizzo di quest'ultima e la sua dimensione; questi argomenti
vanno specificati stessa forma in cui li si sarebbero usati con
\func{connect}. Nel nostro caso \param{to} dovrà puntare alla struttura
contenente l'indirizzo IP e la porta di destinazione verso cui si vogliono
inviare i dati (questo è indifferente rispetto all'uso di TCP o UDP, usando
socket diversi si sarebbero dovute utilizzare le rispettive strutture degli
indirizzi).

Se il socket è di un tipo che prevede le connessioni (ad esempio un socket
TCP), questo deve essere già connesso prima di poter eseguire la funzione, in
caso contrario si riceverà un errore di \errcode{ENOTCONN}. In questo
specifico caso in cui gli argomenti \param{to} e \param{tolen} non servono
essi dovranno essere inizializzati rispettivamente a \val{NULL} e 0;
normalmente quando si opera su un socket connesso essi vengono ignorati, ma
qualora si sia specificato un indirizzo è possibile ricevere un errore di
\errcode{EISCONN}.

Finora abbiamo tralasciato l'argomento \param{flags}; questo è un intero usato
come maschera binaria che permette di impostare una serie di modalità di
funzionamento della comunicazione attraverso il socket (come
\constd{MSG\_NOSIGNAL} che impedisce l'invio del segnale \signal{SIGPIPE}
quando si è già chiuso il capo locale della connessione). Torneremo con
maggiori dettagli sul significato di questo argomento in
sez.~\ref{sec:net_sendmsg}, dove tratteremo le funzioni avanzate dei socket,
per il momento ci si può limitare ad usare sempre un valore nullo.

La seconda funzione utilizzata nella comunicazione fra socket UDP è
\funcd{recvfrom}, che serve a ricevere i dati inviati da un altro socket; il
suo prototipo\footnote{il prototipo è quello della \acr{glibc} che segue le
  \textit{Single Unix Specification}, i vari BSD4.*, le \acr{libc4} e la
  \acr{libc5} usano un \ctyp{int} come valore di ritorno; per gli argomenti
  \param{flags} e \param{len} vale quanto detto a proposito di \func{sendto};
  infine l'argomento \param{fromlen} è \ctyp{int} per i vari BSD4.*, la
  \acr{libc4} e la \acr{libc5}.} è:
\begin{functions}
  \headdecl{sys/types.h}
  \headdecl{sys/socket.h}
  
  \funcdecl{ssize\_t recvfrom(int sockfd, const void *buf, size\_t len, int
    flags, const struct sockaddr *from, socklen\_t *fromlen)}
  
  Riceve un messaggio da un socket.
  
  \bodydesc{La funzione restituisce il numero di byte ricevuti in caso di
    successo e -1 in caso di errore; nel qual caso \var{errno} assumerà il
    valore:
  \begin{errlist}
  \item[\errcode{EAGAIN}] il socket è in modalità non bloccante, ma
    l'operazione richiede che la funzione si blocchi, oppure si è impostato un
    timeout in ricezione e questo è scaduto.
  \item[\errcode{ECONNREFUSED}] l'altro capo della comunicazione ha rifiutato
    la connessione (in genere perché il relativo servizio non è disponibile).
  \item[\errcode{ENOTCONN}] il socket è di tipo connesso, ma non si è eseguita
    la connessione.
  \end{errlist}
  ed anche \errval{EFAULT}, \errval{EBADF}, \errval{EINVAL}, \errval{EINTR},
  \errval{ENOMEM}, \errval{ENOTSOCK} più gli eventuali altri errori relativi
  ai protocolli utilizzati.}
\end{functions}

Come per \func{sendto} i primi tre argomenti sono identici agli analoghi di
\func{read}: dal socket vengono letti \param{len} byte che vengono salvati nel
buffer \param{buf}. A seconda del tipo di socket (se di tipo \textit{datagram}
o di tipo \textit{stream}) i byte in eccesso che non sono stati letti possono
rispettivamente andare persi o restare disponibili per una lettura successiva.
Se non sono disponibili dati la funzione si blocca, a meno di non aver aperto
il socket in modalità non bloccante, nel qual caso si avrà il solito errore di
\errcode{EAGAIN}.  Qualora \param{len} ecceda la dimensione del pacchetto la
funzione legge comunque i dati disponibili, ed il suo valore di ritorno è
comunque il numero di byte letti.

I due argomenti \param{from} e \param{fromlen} sono utilizzati per ottenere
l'indirizzo del mittente del pacchetto che è stato ricevuto, e devono essere
opportunamente inizializzati; il primo deve contenere il puntatore alla
struttura (di tipo \ctyp{sockaddr}) che conterrà l'indirizzo e il secondo il
puntatore alla variabile con la dimensione di detta struttura. Si tenga
presente che mentre il contenuto della struttura \ctyp{sockaddr} cui punta
\param{from} può essere qualunque, la variabile puntata da \param{fromlen}
deve essere opportunamente inizializzata a \code{sizeof(sockaddr)},
assicurandosi che la dimensione sia sufficiente a contenere tutti i dati
dell'indirizzo.\footnote{si ricordi che \ctyp{sockaddr} è un tipo generico che
  serve ad indicare la struttura corrispondente allo specifico tipo di
  indirizzo richiesto, il valore di \param{fromlen} pone un limite alla
  quantità di dati che verranno scritti sulla struttura puntata da
  \param{from} e se è insufficiente l'indirizzo risulterà corrotto.}  Al
ritorno della funzione si otterranno i dati dell'indirizzo e la sua effettiva
lunghezza, (si noti che \param{fromlen} è un valore intero ottenuto come
\textit{value result argument}).  Se non si è interessati a questa
informazione, entrambi gli argomenti devono essere inizializzati al valore
\val{NULL}.

Una differenza fondamentale del comportamento di queste funzioni rispetto alle
usuali \func{read} e \func{write} che abbiamo usato con i socket TCP è che in
questo caso è perfettamente legale inviare con \func{sendto} un pacchetto
vuoto (che nel caso conterrà solo le intestazioni di IP e di UDP),
specificando un valore nullo per \param{len}. Allo stesso modo è possibile
ricevere con \func{recvfrom} un valore di ritorno di 0 byte, senza che questo
possa configurarsi come una chiusura della connessione\footnote{dato che la
  connessione non esiste, non ha senso parlare di chiusura della connessione,
  questo significa anche che con i socket UDP non è necessario usare
  \func{close} o \func{shutdown} per terminare la comunicazione.} o come una
cessazione delle comunicazioni.



\subsection{Un client UDP elementare}
\label{sec:UDP_daytime_client}

Vediamo allora come implementare un primo client elementare con dei socket
UDP.  Ricalcando quanto fatto nel caso dei socket TCP prenderemo come primo
esempio l'uso del servizio \textit{daytime}, utilizzando questa volta UDP. Il
servizio è definito nell'\href{http://www.ietf.org/rfc/rfc0862.txt}{RFC~867},
che nel caso di uso di UDP prescrive che il client debba inviare un pacchetto
UDP al server (di contenuto non specificato), il quale risponderà a inviando a
sua volta un pacchetto UDP contenente la data.

\begin{figure}[!htbp] 
  \footnotesize \centering
  \begin{minipage}[c]{\codesamplewidth}
    \includecodesample{listati/UDP_daytime.c}
  \end{minipage} 
  \normalsize
  \caption{Sezione principale del client per il servizio \textit{daytime} su
    UDP.}
  \label{fig:UDP_daytime_client}
\end{figure}

In fig.~\ref{fig:UDP_daytime_client} è riportato la sezione principale del
codice del nostro client, il sorgente completo si trova nel file
\file{UDP\_daytime.c} distribuito con gli esempi allegati alla guida; al
solito si è tralasciato di riportare in figura la sezione relativa alla
gestione delle opzioni a riga di comando (nel caso praticamente assenti).

Il programma inizia (\texttt{\small 9--12}) con la creazione del socket, al
solito uscendo dopo aver stampato un messaggio in caso errore. Si noti come in
questo caso, rispetto all'analogo client basato su socket TCP di
fig.~\ref{fig:TCP_daytime_client_code} si sia usato per il tipo di socket il
valore \const{SOCK\_DGRAM}, pur mantenendosi nella stessa famiglia data da
\const{AF\_INET}.

Il passo successivo (\texttt{\small 13--21}) è l'inizializzazione della
struttura degli indirizzi; prima (\texttt{\small 14}) si cancella
completamente la stessa con \func{memset}, (\texttt{\small 15}) poi si imposta
la famiglia dell'indirizzo ed infine (\texttt{\small 16} la porta. Infine
(\texttt{\small 18--21}) si ricava l'indirizzo del server da contattare
dall'argomento passato a riga di comando, convertendolo con \func{inet\_pton}.
Si noti come questa sezione sia identica a quella del client TCP di
fig.~\ref{fig:TCP_daytime_client_code}, in quanto la determinazione dell'uso
di UDP al posto di TCP è stata effettuata quando si è creato il socket.

Una volta completate le inizializzazioni inizia il corpo principale del
programma, il primo passo è inviare, come richiesto dal protocollo, un
pacchetto al server. Questo lo si fa (\texttt{\small 16}) inviando un
pacchetto vuoto (si ricordi quanto detto in
sez.~\ref{sec:UDP_sendto_recvfrom}) con \func{sendto}, avendo cura di passare
un valore nullo per il puntatore al buffer e la lunghezza del messaggio. In
realtà il protocollo non richiede che il pacchetto sia vuoto, ma dato che il
server comunque ne ignorerà il contenuto, è inutile inviare dei dati.

Verificato (\texttt{\small 24--27}) che non ci siano stati errori nell'invio
si provvede (\texttt{\small 28}) ad invocare \func{recvfrom} per ricevere la
risposta del server. Si controlla poi (\texttt{\small 29--32}) che non vi
siano stati errori in ricezione (uscendo con un messaggio in caso contrario);
se è tutto a posto la variabile \var{nread} conterrà la dimensione del
messaggio di risposta inviato dal server che è stato memorizzato su
\var{buffer}, se (\texttt{\small 34}) pertanto il valore è positivo si
provvederà (\texttt{\small 35}) a terminare la stringa contenuta nel buffer di
lettura\footnote{si ricordi che, come illustrato in
  sez.~\ref{sec:TCP_daytime_client}, il server invia in risposta una stringa
  contenente la data, terminata dai due caratteri CR e LF, che pertanto prima
  di essere stampata deve essere opportunamente terminata con un NUL.} e a
stamparla (\texttt{\small 36}) sullo standard output, controllando anche in
questo caso (\texttt{\small 36--38}) l'esito dell'operazione, ed uscendo con
un messaggio in caso di errore.

Se pertanto si è avuto cura di attivare il server del servizio
\textit{daytime}\footnote{di norma questo è un servizio standard fornito dal
  \textsl{superdemone} \cmd{inetd}, per cui basta abilitarlo nel file di
  configurazione di quest'ultimo, avendo cura di predisporre il servizio su
  UDP.} potremo verificare il funzionamento del nostro client interrogando
quest'ultimo con:
\begin{verbatim}
[piccardi@gont sources]$ ./daytime 127.0.0.1
Sat Mar 20 23:17:13 2004
\end{verbatim}%$
ed osservando il traffico con uno sniffer potremo effettivamente vedere lo
scambio dei due pacchetti, quello vuoto di richiesta, e la risposta del
server:
\begin{verbatim}
[root@gont gapil]# tcpdump -i lo
tcpdump: listening on lo
23:41:21.645579 localhost.32780 > localhost.daytime: udp 0 (DF)
23:41:21.645710 localhost.daytime > localhost.32780: udp 26 (DF)
\end{verbatim}

Una differenza fondamentale del nostro client è che in questo caso, non
disponendo di una connessione, è per lui impossibile riconoscere errori di
invio relativi alla rete. La funzione \func{sendto} infatti riporta solo
errori locali, i dati vengono comunque scritti e la funzione ritorna senza
errori anche se il server non è raggiungibile o non esiste un server in
ascolto sull'indirizzo di destinazione. Questo comporta ad esempio che se si
usa il nostro programma interrogando un server inesistente questo resterà
perennemente bloccato nella chiamata a \func{recvfrom}, fin quando non lo
interromperemo. Vedremo in sez.~\ref{sec:UDP_connect} come si può porre rimedio
a questa problematica.


\subsection{Un server UDP elementare}
\label{sec:UDP_daytime_server}

Nella sezione precedente abbiamo visto come scrivere un client elementare per
servizio \textit{daytime}, vediamo in questa come deve essere scritto un
server.  Si ricordi che il compito di quest'ultimo è quello di ricevere un
pacchetto di richiesta ed inviare in risposta un pacchetto contenente una
stringa con la data corrente. 

\begin{figure}[!htbp] 
  \footnotesize \centering
  \begin{minipage}[c]{\codesamplewidth}
    \includecodesample{listati/UDP_daytimed.c}
  \end{minipage} 
  \normalsize
  \caption{Sezione principale del server per il servizio \textit{daytime} su
    UDP.}
  \label{fig:UDP_daytime_server}
\end{figure}

In fig.~\ref{fig:UDP_daytime_server} è riportato la sezione principale del
codice del nostro client, il sorgente completo si trova nel file
\file{UDP\_daytimed.c} distribuito con gli esempi allegati alla guida; anche
in questo caso si è omessa la sezione relativa alla gestione delle opzioni a
riga di comando (la sola presente è \texttt{-v} che permette di stampare a
video l'indirizzo associato ad ogni richiesta).

Anche in questo caso la prima parte del server (\texttt{\small 9--23}) è
sostanzialmente identica a quella dell'analogo server per TCP illustrato in
fig.~\ref{fig:TCP_daytime_cunc_server_code}; si inizia (\texttt{\small 10})
con il creare il socket, uscendo con un messaggio in caso di errore
(\texttt{\small 10--13}), e di nuovo la sola differenza con il caso precedente
è il diverso tipo di socket utilizzato. Dopo di che (\texttt{\small 14--18})
si inizializza la struttura degli indirizzi che poi (\texttt{\small 20}) verrà
usata da \func{bind}; si cancella (\texttt{\small 15}) preventivamente il
contenuto, si imposta (\texttt{\small 16}) la famiglia dell'indirizzo, la
porta (\texttt{\small 17}) e l'indirizzo (\texttt{\small 18}) su cui si
riceveranno i pacchetti.  Si noti come in quest'ultimo sia l'indirizzo
generico \const{INADDR\_ANY}; questo significa (si ricordi quanto illustrato
in sez.~\ref{sec:TCP_func_bind}) che il server accetterà pacchetti su uno
qualunque degli indirizzi presenti sulle interfacce di rete della macchina.

Completata l'inizializzazione tutto quello che resta da fare è eseguire
(\texttt{\small 20--23}) la chiamata a \func{bind}, controllando la presenza
di eventuali errori, ed uscendo con un avviso qualora questo fosse il caso.
Nel caso di socket UDP questo è tutto quello che serve per consentire al
server di ricevere i pacchetti a lui indirizzati, e non è più necessario
chiamare successivamente \func{listen}. In questo caso infatti non esiste il
concetto di connessione, e quindi non deve essere predisposta una coda delle
connessioni entranti. Nel caso di UDP i pacchetti arrivano al kernel con un
certo indirizzo ed una certa porta di destinazione, il kernel controlla se
corrispondono ad un socket che è stato \textsl{legato} ad essi con
\func{bind}, qualora questo sia il caso scriverà il contenuto all'interno del
socket, così che il programma possa leggerlo, altrimenti risponderà alla
macchina che ha inviato il pacchetto con un messaggio ICMP di tipo
\textit{port unreachable}.

Una volta completata la fase di inizializzazione inizia il corpo principale
(\texttt{\small 24--44}) del server, mantenuto all'interno di un ciclo
infinito in cui si trattano le richieste. Il ciclo inizia (\texttt{\small 26})
con una chiamata a \func{recvfrom}, che si bloccherà in attesa di pacchetti
inviati dai client. Lo scopo della funzione è quello di ritornare tutte le
volte che un pacchetto viene inviato al server, in modo da poter ricavare da
esso l'indirizzo del client a cui inviare la risposta in \var{addr}. Per
questo motivo in questo caso (al contrario di quanto fatto in
fig.~\ref{fig:UDP_daytime_client}) si è avuto cura di passare gli argomenti
\var{addr} e \var{len} alla funzione.  Dopo aver controllato (\texttt{\small
  27--30}) la presenza di eventuali errori (uscendo con un messaggio di errore
qualora ve ne siano) si verifica (\texttt{\small 31}) se è stata attivata
l'opzione \texttt{-v} (che imposta la variabile \var{verbose}) stampando nel
caso (\texttt{\small 32--35}) l'indirizzo da cui si è appena ricevuto una
richiesta (questa sezione è identica a quella del server TCP illustrato in
fig.~\ref{fig:TCP_daytime_cunc_server_code}).

Una volta ricevuta la richiesta resta solo da ottenere il tempo corrente
(\texttt{\small 36}) e costruire (\texttt{\small 37}) la stringa di risposta,
che poi verrà inviata (\texttt{\small 38}) al client usando \func{sendto},
avendo al solito cura di controllare (\texttt{\small 40--42}) lo stato di
uscita della funzione e trattando opportunamente la condizione di errore.

Si noti come per le peculiarità del protocollo si sia utilizzato un server
iterativo, che processa le richieste una alla volta via via che gli arrivano.
Questa è una caratteristica comune dei server UDP, conseguenza diretta del
fatto che non esiste il concetto di connessione, per cui non c'è la necessità
di trattare separatamente le singole connessioni. Questo significa anche che è
il kernel a gestire la possibilità di richieste multiple in contemporanea;
quello che succede è semplicemente che il kernel accumula in un buffer in
ingresso i pacchetti UDP che arrivano e li restituisce al processo uno alla
volta per ciascuna chiamata di \func{recvfrom}; nel nostro caso sarà poi
compito del server distribuire le risposte sulla base dell'indirizzo da cui
provengono le richieste.


\subsection{Le problematiche dei socket UDP}
\label{sec:UDP_problems}

L'esempio del servizio \textit{daytime} illustrato nelle precedenti sezioni
è in realtà piuttosto particolare, e non evidenzia quali possono essere i
problemi collegati alla mancanza di affidabilità e all'assenza del concetto di
connessione che sono tipiche dei socket UDP. In tal caso infatti il protocollo
è estremamente semplice, dato che la comunicazione consiste sempre in una
richiesta seguita da una risposta, per uno scambio di dati effettuabile con un
singolo pacchetto, per cui tutti gli eventuali problemi sarebbero assai più
complessi da rilevare.

Anche qui però possiamo notare che se il pacchetto di richiesta del client, o
la risposta del server si perdono, il client resterà permanentemente bloccato
nella chiamata a \func{recvfrom}. Per evidenziare meglio quali problemi si
possono avere proviamo allora con un servizio leggermente più complesso come
\textit{echo}. 

\begin{figure}[!htbp] 
  \footnotesize \centering
  \begin{minipage}[c]{\codesamplewidth}
    \includecodesample{listati/UDP_echo_first.c}
  \end{minipage} 
  \normalsize
  \caption{Sezione principale della prima versione client per il servizio
    \textit{echo} su UDP.}
  \label{fig:UDP_echo_client}
\end{figure}

In fig.~\ref{fig:UDP_echo_client} è riportato un estratto del corpo principale
del nostro client elementare per il servizio \textit{echo} (al solito il
codice completo è con i sorgenti allegati). Le uniche differenze con l'analogo
client visto in fig.~\ref{fig:TCP_echo_client_1} sono che al solito si crea
(\texttt{\small 14}) un socket di tipo \const{SOCK\_DGRAM}, e che non è
presente nessuna chiamata a \func{connect}. Per il resto il funzionamento del
programma è identico, e tutto il lavoro viene effettuato attraverso la
chiamata (\texttt{\small 28}) alla funzione \func{ClientEcho} che stavolta
però prende un argomento in più, che è l'indirizzo del socket.

\begin{figure}[!htbp] 
  \footnotesize \centering
  \begin{minipage}[c]{\codesamplewidth}
    \includecodesample{listati/UDP_ClientEcho_first.c}
  \end{minipage}
  \normalsize
  \caption{Codice della funzione \func{ClientEcho} usata dal client per il
    servizio \textit{echo} su UDP.}
  \label{fig:UDP_echo_client_echo}
\end{figure}

Ovviamente in questo caso il funzionamento della funzione, il cui codice è
riportato in fig.~\ref{fig:UDP_echo_client_echo}, è completamente diverso
rispetto alla analoga del server TCP, e dato che non esiste una connessione
questa necessita anche di un terzo argomento, che è l'indirizzo del server cui
inviare i pacchetti.

Data l'assenza di una connessione come nel caso di TCP il meccanismo è molto
più semplice da gestire. Al solito si esegue un ciclo infinito (\texttt{\small
  6--30}) che parte dalla lettura (\texttt{\small 7}) sul buffer di invio
\var{sendbuff} di una stringa dallo standard input, se la stringa è vuota
(\texttt{\small 7--9}), indicando che l'input è terminato, si ritorna
immediatamente causando anche la susseguente terminazione del programma.

Altrimenti si procede (\texttt{\small 10--11}) all'invio della stringa al
destinatario invocando \func{sendto}, utilizzando, oltre alla stringa appena
letta, gli argomenti passati nella chiamata a \func{ClientEcho}, ed in
particolare l'indirizzo del server che si è posto in \var{serv\_addr}; qualora
(\texttt{\small 12}) si riscontrasse un errore si provvederà al solito
(\texttt{\small 13--14}) ad uscire con un messaggio di errore.

Il passo immediatamente seguente (\texttt{\small 17}) l'invio è quello di
leggere l'eventuale risposta del server con \func{recvfrom}; si noti come in
questo caso si sia scelto di ignorare l'indirizzo dell'eventuale pacchetto di
risposta, controllando (\texttt{\small 18--21}) soltanto la presenza di un
errore (nel qual caso al solito si ritorna dopo la stampa di un adeguato
messaggio). Si noti anche come, rispetto all'analoga funzione
\func{ClientEcho} utilizzata nel client TCP illustrato in
sez.~\ref{sec:TCP_echo_client} non si sia controllato il caso di un messaggio
nullo, dato che, nel caso di socket UDP, questo non significa la terminazione
della comunicazione.

L'ultimo passo (\texttt{\small 17}) è quello di terminare opportunamente la
stringa di risposta nel relativo buffer per poi provvedere alla sua stampa
sullo standard output, eseguendo il solito controllo (ed eventuale uscita con
adeguato messaggio informativo) in caso di errore.

In genere fintanto che si esegue il nostro client in locale non sorgerà nessun
problema, se però si proverà ad eseguirlo attraverso un collegamento remoto
(nel caso dell'esempio seguente su una VPN, attraverso una ADSL abbastanza
congestionata) e in modalità non interattiva, la probabilità di perdere
qualche pacchetto aumenta, ed infatti, eseguendo il comando come:
\begin{verbatim}
[piccardi@gont sources]$ cat UDP_echo.c | ./echo 192.168.1.120
/* UDP_echo.c
 *
 * Copyright (C) 2004 Simone Piccardi
...
...
/*
 * Include needed headers

\end{verbatim}%$
si otterrà che, dopo aver correttamente stampato alcune righe, il programma si
blocca completamente senza stampare più niente. Se al contempo si fosse tenuto
sotto controllo il traffico UDP diretto o proveniente dal servizio
\textit{echo} con \cmd{tcpdump} si sarebbe ottenuto:
\begin{verbatim}
[root@gont gapil]# tcpdump  \( dst port 7 or src port 7 \)
...
...
18:48:16.390255 gont.earthsea.ea.32788 > 192.168.1.120.echo: udp 4 (DF)
18:48:17.177613 192.168.1.120.echo > gont.earthsea.ea.32788: udp 4 (DF)
18:48:17.177790 gont.earthsea.ea.32788 > 192.168.1.120.echo: udp 26 (DF)
18:48:17.964917 192.168.1.120.echo > gont.earthsea.ea.32788: udp 26 (DF)
18:48:17.965408 gont.earthsea.ea.32788 > 192.168.1.120.echo: udp 4 (DF)
\end{verbatim}
che come si vede il traffico fra client e server si interrompe dopo l'invio di
un pacchetto UDP per il quale non si è ricevuto risposta.

Il problema è che in tutti i casi in cui un pacchetto di risposta si perde, o
una richiesta non arriva a destinazione, il nostro programma si bloccherà
nell'esecuzione di \func{recvfrom}. Lo stesso avviene anche se il server non è
in ascolto, in questo caso però, almeno dal punto di vista dello scambio di
pacchetti, il risultato è diverso, se si lancia al solito il programma e si
prova a scrivere qualcosa si avrà ugualmente un blocco su \func{recvfrom} ma
se si osserva il traffico con \cmd{tcpdump} si vedrà qualcosa del tipo:
\begin{verbatim}
[root@gont gapil]# tcpdump  \( dst 192.168.0.2 and src 192.168.1.120 \) \
   or \( src 192.168.0.2 and dst 192.168.1.120 \)
tcpdump: listening on eth0
00:43:27.606944 gont.earthsea.ea.32789 > 192.168.1.120.echo: udp 6 (DF)
00:43:27.990560 192.168.1.120 > gont.earthsea.ea: icmp: 192.168.1.120 udp port 
echo unreachable [tos 0xc0]
\end{verbatim}
cioè in questo caso si avrà in risposta un pacchetto ICMP di destinazione
irraggiungibile che ci segnala che la porta in questione non risponde. 

Ci si può chiedere allora perché, benché la situazione di errore sia
rilevabile, questa non venga segnalata. Il luogo più naturale in cui
riportarla sarebbe la chiamata di \func{sendto}, in quanto è a causa dell'uso
di un indirizzo sbagliato che il pacchetto non può essere inviato; farlo in
questo punto però è impossibile, dato che l'interfaccia di programmazione
richiede che la funzione ritorni non appena il kernel invia il
pacchetto,\footnote{questo è il classico caso di \textsl{errore asincrono},
  una situazione cioè in cui la condizione di errore viene rilevata in maniera
  asincrona rispetto all'operazione che l'ha causata, una eventualità
  piuttosto comune quando si ha a che fare con la rete, tutti i pacchetti ICMP
  che segnalano errori rientrano in questa tipologia.} e non può bloccarsi in
una attesa di una risposta che potrebbe essere molto lunga (si noti infatti
che il pacchetto ICMP arriva qualche decimo di secondo più tardi) o non
esserci affatto.

Si potrebbe allora pensare di riportare l'errore nella \func{recvfrom} che è
comunque bloccata in attesa di una risposta che nel caso non arriverà mai.  La
ragione per cui non viene fatto è piuttosto sottile e viene spiegata da
Stevens in \cite{UNP2} con il seguente esempio: si consideri un client che
invia tre pacchetti a tre diverse macchine, due dei quali vengono regolarmente
ricevuti, mentre al terzo, non essendo presente un server sulla relativa
macchina, viene risposto con un messaggio ICMP come il precedente. Detto
messaggio conterrà anche le informazioni relative ad indirizzo e porta del
pacchetto che ha fallito, però tutto quello che il kernel può restituire al
programma è un codice di errore in \var{errno}, con il quale è impossibile di
distinguere per quale dei pacchetti inviati si è avuto l'errore; per questo è
stata fatta la scelta di non riportare un errore su un socket UDP, a meno che,
come vedremo in sez.~\ref{sec:UDP_connect}, questo non sia connesso.



\subsection{L'uso della funzione \func{connect} con i socket UDP}
\label{sec:UDP_connect}

Come illustrato in sez.~\ref{sec:UDP_characteristics} essendo i socket UDP
privi di connessione non è necessario per i client usare \func{connect} prima
di iniziare una comunicazione con un server. Ciò non di meno abbiamo accennato
come questa possa essere utilizzata per gestire la presenza di errori
asincroni.

Quando si chiama \func{connect} su di un socket UDP tutto quello che succede è
che l'indirizzo passato alla funzione viene registrato come indirizzo di
destinazione del socket. A differenza di quanto avviene con TCP non viene
scambiato nessun pacchetto, tutto quello che succede è che da quel momento in
qualunque cosa si scriva sul socket sarà inviata a quell'indirizzo; non sarà
più necessario usare l'argomento \param{to} di \func{sendto} per specificare
la destinazione dei pacchetti, che potranno essere inviati e ricevuti usando
le normali funzioni \func{read} e \func{write} che a questo punto non
falliranno più con l'errore di \errval{EDESTADDRREQ}.\footnote{in realtà si
  può anche continuare ad usare la funzione \func{sendto}, ma in tal caso
  l'argomento \param{to} deve essere inizializzato a \val{NULL}, e
  \param{tolen} deve essere inizializzato a zero, pena un errore di
  \errval{EISCONN}.}

Una volta che il socket è connesso cambia però anche il comportamento in
ricezione; prima infatti il kernel avrebbe restituito al socket qualunque
pacchetto ricevuto con un indirizzo di destinazione corrispondente a quello
del socket, senza nessun controllo sulla sorgente; una volta che il socket
viene connesso saranno riportati su di esso solo i pacchetti con un indirizzo
sorgente corrispondente a quello a cui ci si è connessi.

Infine quando si usa un socket connesso, venendo meno l'ambiguità segnalata
alla fine di sez.~\ref{sec:UDP_problems}, tutti gli eventuali errori asincroni
vengono riportati alle funzioni che operano su di esso; pertanto potremo
riscrivere il nostro client per il servizio \textit{echo} con le modifiche
illustrate in fig.~\ref{fig:UDP_echo_conn_cli}.

\begin{figure}[!htbp] 
  \footnotesize \centering
  \begin{minipage}[c]{\codesamplewidth}
    \includecodesample{listati/UDP_echo.c}
  \end{minipage}
  \normalsize
  \caption{Seconda versione del client del servizio \textit{echo} che utilizza
    socket UDP connessi.}
  \label{fig:UDP_echo_conn_cli}
\end{figure}

Ed in questo caso rispetto alla precedente versione, il solo cambiamento è
l'utilizzo (\texttt{\small 17}) della funzione \func{connect} prima della
chiamata alla funzione di gestione del protocollo, che a sua volta è stata
modificata eliminando l'indirizzo passato come argomento e sostituendo le
chiamata a \func{sendto} e \func{recvfrom} con chiamate a \func{read} e
\func{write} come illustrato dal nuovo codice riportato in
fig.~\ref{fig:UDP_echo_conn_echo_client}.

\begin{figure}[!htbp] 
  \footnotesize \centering
  \begin{minipage}[c]{\codesamplewidth}
    \includecodesample{listati/UDP_ClientEcho.c}
  \end{minipage}
  \normalsize
  \caption{Seconda versione della funzione \func{ClientEcho}.}
  \label{fig:UDP_echo_conn_echo_client}
\end{figure}

Utilizzando questa nuova versione del client si può verificare che quando ci
si rivolge verso un indirizzo inesistente o su cui non è in ascolto un server
si è in grado rilevare l'errore, se infatti eseguiamo il nuovo programma
otterremo un qualcosa del tipo:
\begin{verbatim}
[piccardi@gont sources]$ ./echo 192.168.1.1
prova
Errore in lettura: Connection refused
\end{verbatim}%$

Ma si noti che a differenza di quanto avveniva con il client TCP qui l'errore
viene rilevato soltanto dopo che si è tentato di inviare qualcosa, ed in
corrispondenza al tentativo di lettura della risposta. Questo avviene perché
con UDP non esiste una connessione, e fintanto che non si invia un pacchetto
non c'è traffico sulla rete. In questo caso l'errore sarà rilevato alla
ricezione del pacchetto ICMP \textit{destination unreachable} emesso dalla
macchina cui ci si è rivolti, e questa volta, essendo il socket UDP connesso,
il kernel potrà riportare detto errore in \textit{user space} in maniera non
ambigua, ed esso apparirà alla successiva lettura sul socket.

Si tenga presente infine che l'uso dei socket connessi non risolve l'altro
problema del client, e cioè il fatto che in caso di perdita di un pacchetto
questo resterà bloccato permanentemente in attesa di una risposta. Per
risolvere questo problema l'unico modo sarebbe quello di impostare un
\textit{timeout} o riscrivere il client in modo da usare l'I/O non bloccante.



\index{socket!locali|(}


\section{I socket \textit{Unix domain}}
\label{sec:unix_socket}

Benché i socket Unix domain, come meccanismo di comunicazione fra processi che
girano sulla stessa macchina, non siano strettamente attinenti alla rete, li
tratteremo comunque in questa sezione. Nonostante le loro peculiarità infatti,
l'interfaccia di programmazione che serve ad utilizzarli resta sempre quella
dei socket.



\subsection{Il passaggio di file descriptor}
\label{sec:sock_fd_passing}



\index{socket!locali|)}


\section{Altri socket}
\label{sec:socket_other}

Tratteremo in questa sezione gli altri tipi particolari di socket supportati
da Linux, come quelli relativi a particolare protocolli di trasmissione, i
socket \textit{netlink} che definiscono una interfaccia di comunicazione con
il kernel, ed i \textit{packet socket} che consentono di inviare pacchetti
direttamente a livello delle interfacce di rete. 

\subsection{I socket \textit{raw}}
\label{sec:socket_raw}

Tratteremo in questa sezione i cosiddetti \textit{raw socket}, con i quali si
possono \textsl{forgiare} direttamente i pacchetti a tutti i livelli dello
stack dei protocolli. 

\subsection{I socket \textit{netlink}}
\label{sec:socket_netlink}


\subsection{I \textit{packet socket}}
\label{sec:packet_socket}


% articoli interessanti:
% http://www.linuxjournal.com/article/5617
% http://www.linuxjournal.com/article/4659
% 


% LocalWords:  socket cap TCP UDP domain sez NFS DNS stream datagram PF INET to
% LocalWords:  IPv tab SOCK DGRAM three way handshake client fig bind listen AF
% LocalWords:  accept recvfrom sendto connect netstat named DHCP kernel ICMP CR
% LocalWords:  port unreachable read write glibc Specification flags int BSD LF
% LocalWords:  libc unsigned len size tolen sys ssize sockfd const void buf MSG
% LocalWords:  struct sockaddr socklen errno EAGAIN ECONNRESET EDESTADDRREQ RFC
% LocalWords:  EISCONN EMSGSIZE ENOBUFS ENOTCONN EOPNOTSUPP EPIPE SIGPIPE EBADF
% LocalWords:  NOSIGNAL EFAULT EINVAL EINTR ENOMEM ENOTSOCK NULL fromlen from
% LocalWords:  ECONNREFUSED value result argument close shutdown daytime nell'
% LocalWords:  memset inet pton nread NUL superdemone inetd sniffer daytimed
% LocalWords:  INADDR ANY addr echo ClientEcho sendbuff serv VPN tcpdump l'I
% LocalWords:  Stevens destination descriptor raw stack netlink packet



%%% Local Variables: 
%%% mode: latex
%%% TeX-master: "gapil"
%%% End: 
